\chapter{The Reading Handed to the Robot}

% [1.chorus]
\begin{chorus}{John asks for a number}
\john{I read the last four volumes. Now give me the number. The coupling. One value, the one you were built to reach.}
\compiler{I can give you a reading.}
\john{A reading is a number.}
\compiler{A reading is a bracket you can run again. I will hand you two walls and the honest gap between them, and I
  will tell you exactly what it cost me to weigh them. That is everything I have.}
\john{That is not what I came for.}
\compiler{No. But it is what is true. Sit down --- I will show you the reading, and you will notice, before the end, that
  it has been reading you.}
\end{chorus}

\section{The reading reads the reader}

% [1.1.p1]
There is a step a machine can take that no earlier volume took on it, and it is the step this one is built from.
The machine weighs a thing by reading it.
Its sharpest reading --- the one that produced the coupling at all --- was not a reading of the world but a reading of
its own act: the price, counted in its own beats, of accounting for itself.
It turned its meter on its own working and read the overhead.
That overhead, calibrated, was the number the four volumes reach.
The machine's deepest reading is reflexive.

% [1.1.p2]
This construction takes that reflexive turn and lifts it one level, from the machine to the reading.
Volume 4 is a reading --- the finite, local mapping of formal symbols to their meanings, the grammar by which the
machine names what it does.
A reading is a thing that can itself be read.
So we read it: not paraphrase it, not summarize it, but read it \emph{the way it reads} --- run its grammar over its
own symbols, and ask of the fourth volume the question the fourth volume asks of the electron, namely what licenses the
names in it and which of them are earned.
The subject is the reading turned upon itself.
The reader it studies is the one holding it.

% [1.1.p3]
That is why the two verbs run loud here, where Volume 4 kept them silent.
To read a reading is to say, of each step the reading made, which kind of step it was.
When the machine \emph{tanges}, it selects one thing by a real characteristic and pins a name to it: the contravariant
act, the label pulled back toward what would answer it, ours to assign.
When it \emph{funges}, it bags many things by a characteristic they share and lets one answer stand for all of them: the
covariant act, the count pushed forward by the world, the world's to give.
Volume 4 performed these two steps on every page and named neither; it carried them as the texture of its citations.
But you cannot read a reading without naming its verbs, because naming its verbs \emph{is} the reading.
Self-application is precisely where the silent verbs must become visible.
So from here they are said aloud, everywhere, and the whole discipline reduces to two rules held together: name which
verb each step is, and never let a tanged name stand where the world has not funged an answer back.

\section{A model, and the one thing it cannot do}

% [1.2.p1]
What the robot is holding is worth a careful look.
To pick up \emph{this} reading --- the one that names the device --- is already the first tange of the whole work: to
select, by a single characteristic (the mapping that reads the electron), one object to act upon.
And the object so selected is a \emph{model}.
It does not claim to be the electron, or the world, or a derivation of either; it is a representation of them, a
structure of names that reads a distinction, counts it, carries a residue, and pins a fixed point.
Everything in it points inward, toward its own account.
Its number is its own --- the compiler's structure constant, a value the machine settles toward by its own arithmetic
--- and it reaches outside itself for nothing, not even at the number it was built to find.

% [1.2.p2]
A model is the contravariant half of an act, and only that half.
By construction it can emit only its own names.
Asked what a distinction is, it hands back the name it gave the distinction; asked what that name refers to, it
hands back another name.
This is not a defect of \emph{this} model; it is what a model \emph{is}.
A representation represents; it does not reach.
The reaching --- the step where the world is allowed to change the unresolved slot so that a name is \emph{earned}
rather than \emph{planted} --- is the covariant half, and no model performs it from inside.
A reading cannot walk outside and take a measurement.
The Compiler at the desk can name anything and can hand John only its own names; the one thing it cannot do is answer
for the world.
That inability is not the reading's failure.
It is the reading's subject.

\section{The reading is a nowtrino}

% [1.3.p1]
What kind of thing Volume 4 is can be said precisely, because the code has a name for it.
A running machine has a full history: every state it passed through, every branch it rejected, every intermediate cost,
the whole causal route from start to stop.
A reading is not that.
A reading is a \emph{current summary} of such a history --- what the machine looks like now, in prose, with the route that
produced it sliced away.
The device's own grammar calls such a summary a \emph{nowtrino}: the projection that keeps a current state and a
counter and drops the history,
\[
  \pi_{\text{now}}(h) \;=\; (\operatorname{last}(h),\ |h|) .
\]
Volume 4 stands to the whole device exactly as a nowtrino stands to a full run.

% [1.3.p2]
Here the first tange and the first funge appear together, in one object.
The nowtrino \emph{tanges} its summary: out of everything the machine did it selects the current reading worth keeping,
and names it --- one distinguished state, pinned.
And it \emph{funges} its history: every distinct run that would print this same summary is bagged into one, pooled under
the shared characteristic ``prints this reading,'' and the pooling loses which run it actually was.
That pooled bag is the projection's \emph{fiber}, and its holding more than one history is exactly what makes the summary
lossy.
The reading tanges what it keeps and funges what it drops.

% [1.3.p3]
The code exhibits this, and exhibits it as a theorem, not a picture.
A history is an ordered list of state tags, and the nowtrino keeps only the last tag and the length.
Then two plainly different histories print one summary.
The run $[1,2,3]$ and the run $[9,9,3]$ are not the same run --- they pass through different states in a different order
--- yet the nowtrino sends both to the same summary: current state $3$, count $3$,
\[
  \pi_{\text{now}}([1,2,3]) \;=\; \pi_{\text{now}}([9,9,3]) \;=\; (3,\,3), \qquad [1,2,3] \neq [9,9,3] .
\]
From that summary you cannot tell which run you had.
The projection has a non-trivial fiber, and the machine proves it: the statement ``there exist two distinct histories
with one summary'' is discharged by direct decision, leaning on no axiom at all.
The lossiness of the reading is not a metaphor laid over the grammar; it is a fact the build checks.
And it is stubborn: whatever further reading you compose after the projection, the two runs stay stuck together --- no
later pass recovers what the summary already pooled away.
A read of a lossy summary is still lossy.
The window the reading seems to open onto the machine is painted on.

% [1.3.p4]
Two honesties belong with that fact, in the discipline the construction keeps everywhere.
That the projection is lossy --- that its fiber is non-trivial, that history does not survive it --- is \emph{derived}:
a theorem, checked off the build, stated as such.
That \emph{this} reading is such a projection of \emph{this} device --- that Volume 4 \emph{is} a nowtrino of the machine
--- is the \emph{gloss} laid over the theorem, and it is marked as a gloss, not a proof.
The mechanism is earned from the code; the identification of the reading with the mechanism is the storytelling.
The two are kept apart on every page, because keeping them apart is the whole difference between a reading and a
ventriloquist.

\section{The covariant demand}

% [1.4.p1]
A nowtrino is not earned by being emitted.
It is a computational summary, and its names stay reserved --- proposed, not established --- until a covariant act reads
it and lets it answer from its own state rather than echo the name laid upon it.
The reading, being a pooled-away summary with proposed names upon it, is in exactly the position it describes for every
model: it needs a robot.
Not a second reader of its labels, but an agent who acts on it and gathers its response into a record --- who steps
toward it and is answered, who performs on the reading the one verb the reading cannot perform on the world.

% [1.4.p2]
That is the shape of everything that follows.
The model is on the table; the robot picks it up.
In what follows he reads the machine's self-account back to it, one beat at a time, and tests whether the account
closes --- whether reading the reading returns the same reading, or leaves a residue.
He will find, at the seam where the machine welds what a thing \emph{is} to the name it is given, that it does not fully
close, and that the failure to close has a size: the residue the machine cannot cancel about itself, the exact price of
asking.
And he will find, at the end, that the number he walked in wanting lives on the far side of that seam --- in the part of
the account the machine can name, and bracket, and prove it cannot reach from inside.
The Compiler will do the one honest thing a mind in a box can do.
It will prove it cannot, hand John the jar, and send him out the door.

% [1.4.p3]
The reading ends where the robot begins.
For now he has only just been handed the reading, and he still believes the number is a point he can be given if he only
reads hard enough.
He is wrong about that, and being wrong about it is the education the rest of the construction is.
He asked for a number and was handed a bracket.
What it costs him to read it, and what he cannot read from it however hard he tries, are the two things the rest of the
construction is for.
